\section*{第 \ref{ch:b1:classifying-biodiversity} 章——生物世界的分类}

\begin{solution}{exo:b1:classifying-biodiversity:1}
\emph{Felis catus}：属 \emph{Felis}，种加词 \emph{catus}；科
Felidae，目 Carnivora，纲 Mammalia，门 Chordata。
\end{solution}

\begin{solution}{exo:b1:classifying-biodiversity:2}
\omterm{def:b1:classifying-biodiversity:characters}{同源}性：由共同祖先遗传而来的相似性（蝙蝠与人类的前肢
\omterm{def:b1:body-plans-tissues:segments}{骨骼}）。\omterm{def:b1:classifying-biodiversity:characters}{同塑}性：由趋同或
逆转而产生的相似性（蝙蝠与鸟类的翅膀）。只有共有衍生的
\omterm{def:b1:classifying-biodiversity:characters}{同源}性才是\omterm{def:b1:classifying-biodiversity:systematics}{演化支}的证据。
\end{solution}

\begin{solution}{exo:b1:classifying-biodiversity:3}
“鱼类”：\omterm{def:b1:classifying-biodiversity:monophyly}{并系}的（不包括从
鱼类演化而来的四足动物）。鸟类：\omterm{def:b1:classifying-biodiversity:monophyly}{单系}的（包括
同一有羽毛祖先的全部后代）。“恒温动物”：\omterm{def:b1:classifying-biodiversity:monophyly}{多系}的（鸟类与哺乳动物
分别独立获得了恒温性）。
\end{solution}

\begin{solution}{exo:b1:classifying-biodiversity:4}
细菌域、古菌域、真核域。古菌的膜脂通过醚键连接，
且不含\omterm{def:b1:carbohydrates:polysaccharide}{肽聚糖}；其 \omterm{def:b1:nucleic-acids:chain}{RNA} 聚合酶和\omterm{def:b1:gene-expression:ribosome}{核糖体}
\omterm{def:b1:proteins:peptide}{蛋白质}与真核生物的相似。
\end{solution}

\begin{solution}{exo:b1:classifying-biodiversity:5}
细菌不发生\omterm{def:b1:nucleic-acids:melting}{杂交}——它们靠分裂繁殖，并跨越本应划分\omterm{def:b1:classifying-biodiversity:species}{物种}的
界线交换\omterm{def:b1:genomes:gene}{基因}——而化石无法检验其是否\omterm{def:b1:nucleic-acids:melting}{杂交}。
细菌依据序列相似性界定（以
\omterm{def:b1:gene-expression:ribosome}{核糖体} \omterm{def:b1:nucleic-acids:chain}{RNA} 或\omterm{def:b1:genomes:genome}{基因组}的阈值为准）；化石则依据形态。
\end{solution}

\begin{solution}{exo:b1:classifying-biodiversity:6}
树 (X,Y),Z：\omterm{def:b1:classifying-biodiversity:characters}{性状} 1、3、4 各变化一次，\omterm{def:b1:classifying-biodiversity:characters}{性状} 2 变化两次，
长度为 5。树 (Y,Z),X：\omterm{def:b1:classifying-biodiversity:characters}{性状} 2 变化一次，\omterm{def:b1:classifying-biodiversity:characters}{性状} 1、3、4 各变化两次，
长度为 7。优选 (X,Y)；\omterm{def:b1:classifying-biodiversity:characters}{性状} 2 为
\omterm{def:b1:classifying-biodiversity:characters}{同塑}性的结果。
\end{solution}

\begin{solution}{exo:b1:classifying-biodiversity:7}
\omterm{def:b1:classifying-biodiversity:characters}{外群}决定每个\omterm{def:b1:classifying-biodiversity:characters}{性状}的极性：其状态被视为祖先状态，
因此只有衍生状态才算作证据。若\omterm{def:b1:classifying-biodiversity:characters}{外群}实际上位于
类群内部，则会使部分衍生状态看似祖先状态，
并使树根置于错误位置，
使\omterm{def:b1:classifying-biodiversity:systematics}{演化支}变成\omterm{def:b1:classifying-biodiversity:monophyly}{并系}类群。
\end{solution}

\begin{solution}{exo:b1:classifying-biodiversity:8}
\omterm{def:b1:gene-expression:ribosome}{核糖体}与\omterm{def:b1:proteins:peptide}{蛋白质}序列显示，动物与真菌互为姐妹
支系，两者都属于后鞭毛类群，其游动\omterm{def:b1:cell-unit-of-life:cell}{细胞}
（精子、壶菌游动孢子）具有单一的后生鞭毛；植物
则属于另一支系，以初级\omterm{def:b1:photosynthesis:chloroplast}{叶绿体}为界定特征。
\end{solution}

\begin{solution}{exo:b1:classifying-biodiversity:9}
若将产甲烷菌归为细菌内的一个界，则意味着它们与其他细菌
的亲缘关系比与真核生物更近；而序列显示三个支系
彼此距离相等，因此原核生物并非一个类群，生命最深层的
分化是三分的。
\end{solution}

\begin{solution}{exo:b1:classifying-biodiversity:10}
每个树种上的专一性甲虫：$0.2\times 1200 = 240$；
仅限该树的冠层昆虫：$240/0.4 = 600$；整棵树：
$600/(2/3) = 900$；乘以 \num{50000} 棵树，共 4500 万。
假设：熏蒸的那棵树具有代表性，专一性确系真实的，
而非抽样误差；甲虫占每棵树上昆虫的 \qty{40}{\%}，
且仅专一于某一棵树的\omterm{def:b1:classifying-biodiversity:species}{物种}不被重复计数——
每项估计可能偏差两倍或更多，后来的估计
已降至 500 万左右。
\end{solution}

\begin{solution}{exo:b1:classifying-biodiversity:11}
\omterm{def:b1:classifying-biodiversity:nomenclature}{模式标本}是该名称所指的对象：活的
\omterm{def:b1:populations:population}{种群}若被判定与\omterm{def:b1:classifying-biodiversity:nomenclature}{模式标本}同种，则属同一\omterm{def:b1:classifying-biodiversity:species}{物种}，
否则为新\omterm{def:b1:classifying-biodiversity:species}{物种}；在后一种情形下，旧名称仍归属\omterm{def:b1:classifying-biodiversity:nomenclature}{模式标本}，
新\omterm{def:b1:populations:population}{种群}需要新的名称。优先权原则
保证：若发现两个名称指同一\omterm{def:b1:classifying-biodiversity:species}{物种}，则
保留较早的那个，使名称不会随每一种新
观点而增多或漂移。
\end{solution}

\begin{solution}{exo:b1:classifying-biodiversity:12}
\omterm{def:b1:classifying-biodiversity:parsimony}{简约性}仅假定趋同与逆转比共同祖先遗传更罕见，
因此树不应因此被迫需要比数据所要求更多的此类事件。
当趋同常见（如\omterm{def:b1:ecosystem-organization:ecosystem}{生境}相似）时该方法失效；当\omterm{def:b1:classifying-biodiversity:characters}{性状}
彼此相关（同一适应被计为多次）时亦然；
当快速演化的支系积累偶然相似而被拉到一起（“长枝吸引”）时，
该方法同样失效。分子序列提供许多独立\omterm{def:b1:classifying-biodiversity:characters}{性状}，
以及可用于检验同一棵树的替换模型；化石提供
衍生状态出现的先后顺序，以及
分支分化的时间；
独立数据之间的一致性即是检验。
\end{solution}

\begin{solution}{pb:b1:classifying-biodiversity:1}
\textbf{1.} \omterm{def:b1:classifying-biodiversity:characters}{性状} 1，颌：除七鳃鳗外所有类群共有，因此它
界定了内群，并确认七鳃鳗为\omterm{def:b1:classifying-biodiversity:characters}{外群}。
\textbf{2.} 1（7 个类群：有颌脊椎动物），2（6：硬骨脊椎动物），3
（5：肺鱼与四足动物），4（4：四足动物），5（3：羊膜动物），8、9
与 10（各 2：蜥蜴+鸽，小鼠+鸽，蜥蜴+鸽），6 与 7
（各 1）。
\textbf{3.} 毛（小鼠）与羽毛（鸽）：独有状态，
即独有衍征，它们可以诊断该类群，但对其
近缘类群毫无说明。
\textbf{4.} \omterm{def:b1:classifying-biodiversity:characters}{性状} 9（恒温性：小鼠+鸽）与 8 及 10（蜥蜴+鸽）相冲突；
冲突发生于这三类羊膜动物之间。
\textbf{5.} 无肺是祖先状态（七鳃鳗、鲨鱼与鳟鱼
皆无肺），有肺为衍生状态；\omterm{def:b1:classifying-biodiversity:characters}{外群}的状态取为
祖先状态。
\textbf{6.} 三种：三者中任何一个都可以作为离群者——
（蜥蜴，鸽）中小鼠在外，（小鼠，鸽）中蜥蜴在外，
（蜥蜴，小鼠）中鸽在外。只有\omterm{def:b1:classifying-biodiversity:characters}{性状} 8、9
和 10 在三者之间存在差异，因此只有它们与该选择有关。
\textbf{7.} （七鳃鳗，（鲨鱼，（鳟鱼，（肺鱼，（青蛙，（蜥蜴，小鼠，
鸽））））））；每个节点由一个\omterm{def:b1:classifying-biodiversity:characters}{性状}标记。
\textbf{8.} 树 A：（小鼠，（蜥蜴，鸽））；树 B：（蜥蜴，（小鼠，
鸽））。
\textbf{9.} 树 A：\omterm{def:b1:classifying-biodiversity:characters}{性状} 1--8 与 10 各变化一次（共 9 次），\omterm{def:b1:classifying-biodiversity:characters}{性状}
9 变化两次：长度 \num{11}。
\textbf{10.} 树 B：\omterm{def:b1:classifying-biodiversity:characters}{性状} 1--7 与 9 各一次（共 8 次），\omterm{def:b1:classifying-biodiversity:characters}{性状} 8 与
10 各两次：长度 \num{12}。
\textbf{11.} 树 A 少一步。
\textbf{12.} 恒温性：它独立出现了两次，一次在哺乳动物支系中，
一次在鸟类支系中——这是趋同，与两者不同的机制相符
（皮毛与汗液相对于羽毛与气囊）。
\textbf{13.} 若不计 8 与 10：树 A 长度为 9（恒温性两次起源），树
B 长度为 8（全部一致）：应优选 B。该结论
取决于所选\omterm{def:b1:classifying-biodiversity:characters}{性状}；\omterm{def:b1:classifying-biodiversity:characters}{性状}较少时，一次\omterm{def:b1:classifying-biodiversity:characters}{同塑}性
即可能左右结果，这正是需要大量\omterm{def:b1:classifying-biodiversity:characters}{性状}，包括分子
\omterm{def:b1:classifying-biodiversity:characters}{性状}的原因。
\textbf{14.} 否：它们的共同祖先也是
四足动物的祖先，因此必须将四类四足动物全部纳入。
\textbf{15.} 是：羊膜卵。
\textbf{16.} 否：鸽从爬行动物内部演化而来（在此树中它是
蜥蜴的姐妹群，在更完整的树中则是鳄类的姐妹群），因此
不含鸟类的爬行动物为\omterm{def:b1:classifying-biodiversity:monophyly}{并系}类群。
\textbf{17.} 否：它们最近的共同祖先是所有羊膜动物的祖先，
该祖先并不具恒温性；\omterm{def:b1:classifying-biodiversity:monophyly}{多系}类群，按一个
趋同\omterm{def:b1:classifying-biodiversity:characters}{性状}归并。
\textbf{18.} 更接近青蛙：它与四足动物共有肺（\omterm{def:b1:classifying-biodiversity:characters}{性状} 3），
而鳟鱼无肺；去掉四足动物的“鱼类”是
\omterm{def:b1:classifying-biodiversity:monophyly}{并系}类群。
\textbf{19.} 纲 Reptilia、纲 Mammalia、纲 Aves：三个
等级相同，尽管蜥蜴与鸽的分歧远比
二者中任一与哺乳动物的分歧晚近；分类阶元属于约定，而
等级相同并不意味着年代相同或差异程度相同。
\textbf{20.} 哺乳动物：每年新增 \num{20} 种，相对 \num{6500} 种为 \qty{0.3}{\%}；
鱼类：每年新增 \num{300} 种，相对 \num{35000} 种为 \qty{0.9}{\%}——哺乳动物
名录已接近完成，而鱼类名录仍在
快速增长。
\textbf{21.} 九成的序列属于从未被培养或命名的生物：
已描述的原核生物仅占即便一升水中
所含原核生物的很小一部分，名录才刚刚开始。
\textbf{22.} 背景：每年 \num{9} 种。千倍时每年 \num{9000}
种——其中大多数是从未被描述的\omterm{def:b1:classifying-biodiversity:species}{物种}。
\textbf{23.} 其\omterm{def:b1:classifying-biodiversity:characters}{性状}、序列及其在树上的位置均
随之丢失；生命之树的一个分支在绘出之前便被抹去，
任何仅由它界定的\omterm{def:b1:classifying-biodiversity:systematics}{演化支}也随之消失。
\textbf{24.} 科规模大、醒目，且几乎都已被描述
（新科现已罕见），因此其数量接近完整；
在研究充分的类群中，\omterm{def:b1:classifying-biodiversity:species}{物种}数、属数与科数之间的
比例较为规则，因此可以将科数外推至
\omterm{def:b1:classifying-biodiversity:species}{物种}数尚不完整的类群。
\textbf{25.} （七鳃鳗，（鲨鱼，（鳟鱼，（肺鱼，（青蛙，（小鼠，（蜥蜴，
鸽））））））），长度 \num{11}。
\end{solution}
